\section{Introduction}
\label{u:sec:introduction}

Fix \(\nu>0\) and a divergence-free datum
\(u_0\in H^\infty(\R^3;\R^3)\).  Smooth-data local theory gives a unique
classical solution on a maximal interval \([0,T_*)\)
\cite{FujitaKato1964,Kato1984}.  Throughout Part~I, \((u,p)\) denotes that
solution.  The problem is to decide whether \(T_*\) can be finite.  On each strict
subinterval \([0,T]\), smoothness already makes every finite Sobolev norm
finite.  That fact alone says nothing about what happens as \(T\uparrow T_*\):
the bound may depend on \(T\) and diverge at the endpoint.  The proof must
instead control the approach to \(T_*\) with constants that do not depend on
the cutoff.

Leray constructed global finite-energy weak solutions \cite{Leray1934}, and
Caffarelli, Kohn, and Nirenberg bounded the size of the possible singular set
for suitable weak solutions \cite{CKN1982}.  Serrin proved regularity under
scale-balanced spacetime integrability hypotheses \cite{Serrin1962};
Escauriaza, Seregin, and \v{S}ver\'ak treated the endpoint
\(L^\infty_tL^3_x\) case \cite{ESS2003}.  Fujita--Kato and Kato supply the
strong local theory used below at the terminal restart
\cite{FujitaKato1964,Kato1984}, and Koch--Tataru place small-data
global well-posedness in the critical space \(\mathrm{BMO}^{-1}\)
\cite{KochTataru2001}.
Tao gives a localization and compactness formulation of the global problem
\cite{Tao2013}.  The regularity criteria show that control of an appropriate
critical quantity excludes singularity.  The argument below derives the
required control from the compatible finite rectangles of the maximal
velocity--pressure history.  Mixed-index exhaustion and the completed
critical-height relationships use those same rectangles and lead to the same
all-orders Sobolev membership.

Pressure is part of that derivative structure.  With the decay normalization,
incompressibility determines it from the current velocity:
\[
 -\Delta p=\partial_i\partial_j(u_i u_j),
 \qquad
 p=R_iR_j(u_i u_j).
\]
Here \(R_i\) are the Riesz transforms, and repeated indices are summed.
Write
\[
 V_n:=\partial_t^nu,
 \qquad
 Q_n:=\partial_t^np.
\]
Differentiating the equation gives
\begin{align}
 Q_n
 &=R_iR_j\sum_{a=0}^{n}\binom naV_{a,i}V_{n-a,j},
 \label{u:eq:intro-pressure-recurrence}\\
 V_{n+1}
 &=\nu\Delta V_n
 -\sum_{a=0}^{n}\binom na(V_a\cdot\nabla)V_{n-a}
 -\nabla Q_n.
 \label{u:eq:intro-velocity-recurrence}
\end{align}
Thus the time derivatives, spatial derivatives, and pressure derivatives are
coordinates of one solution.  A time derivative raises the order of the
viscous term by two spatial derivatives and distributes the transport and
pressure terms among earlier time derivatives.  The resulting estimates must
track temporal depth and spatial depth independently.

The Sobolev consequences used below have a simple form. For any prescribed
spatial derivative order \(k\), a Sobolev order
\(s>k+3/2\) gives a continuous derivative.  Applying this at every finite
order gives the spatial column
\[
 L^2\cap\bigcap_{j\in\Nzero}\dot H^{j+1/2}
 =\bigcap_{m\in\Nzero}H^m\hookrightarrow C^\infty.
\]
The corresponding one-dimensional embedding,
\(\bigcap_r H^r(0,T_*)\hookrightarrow C^\infty(0,T_*)\), applies to
the temporal-height column of
\(H(t)=\tfrac12\norm{\Lambda^{1/2}u(t)}_2^2\).
The viscous contribution to \(H'\) is
\(-\nu\norm{\Lambda^{3/2}u}_2^2\).  Each further time derivative exposes
the next spatial energy together with differentiated transport and pressure
terms.  Subtracting those terms from \(H^{(k)}\) and dividing by
\((-2\nu)^k\) recovers \(\tfrac12\norm{\Lambda^{k+1/2}u}_2^2\), as shown
in Proposition~\ref{u:prop:completed-height-recurrence}.

For finite \(N,M\in\Nzero\), let
\[
 \mathcal R_{N,M}[u_0;T_*]
 :=\bigl((V_n)_{n=0}^{N},(V_{n+1})_{n=0}^{N};(Q_n)_{n=0}^{N}\bigr)
\]
be the corresponding pressure-complete rectangle of the maximal history.
The datum-generated rectangle theorem places these selected rows on the whole
interval \((0,T_*)\), and enlargement preserves every coordinate already
present.  Reading the compatible family by its temporal and spatial indices
gives
\begin{equation}
 u\in\bigcap_{r,m\in\Nzero}H^r(0,T_*;H^m)
 \label{u:eq:intro-all-orders-intersection}
\end{equation}
on a proposed finite maximal interval.  Reading the same family through the
recurrences \eqref{u:eq:intro-pressure-recurrence}--%
\eqref{u:eq:intro-velocity-recurrence} and the completed heights
\(R_k=E_{k+1/2}\) first gives
\(u\in\bigcap_mL^2(0,T_*;H^m)\); the recurrence then recovers every temporal
and pressure row.  Both readings use the same datum-generated compatible
family.

The rectangle and its norm are different objects.  The rectangle retains
\(V_0,\ldots,V_N\), their next time derivatives, and their pressure rows.
The scalar sum of its selected velocity norms on a strict interval is
\begin{equation}
 \mathfrak R_{N,M}(u;T)
 :=\sum_{n=0}^{N}\norm{V_n}_{L^2(0,T;H^{M+1})}^2
 +\sum_{n=0}^{N}\norm{V_{n+1}}_{L^2(0,T;H^{M-1})}^2.
 \label{u:eq:intro-rectangle-norm}
\end{equation}
Each term measures one coordinate of that compatible rectangle.  Its
full-interval membership, written in mixed-row coordinates, gives
\begin{equation}
 \sup_{0<T<T_*}\mathfrak R_{N,M}(u;T)
 \le\sum_{n=0}^N\left(
 \norm{V_n}_{L^2(0,T_*;H^{M+1})}^2
 +\norm{V_{n+1}}_{L^2(0,T_*;H^{M-1})}^2\right)<\infty.
 \label{u:eq:intro-terminal-rectangle}
\end{equation}
Theorem~\ref{u:thm:datum-rectangles} proves the finiteness of this rectangle norm.
The constant may depend on \(N,M\), but is fixed independently of the
cutoff \(T\).  We do not take a supremum over derivative orders.

Larger rectangles literally restrict to the same rows in every smaller
rectangle:
\begin{equation}
 \left.\mathcal R_{N,M}[u_0;T_*]\right|_{(N',M')}
 =\mathcal R_{N',M'}[u_0;T_*]
 \qquad(N'\le N,\ M'\le M).
 \label{u:eq:intro-compatible-intersection}
\end{equation}
Every finite mixed row occurs in a sufficiently large rectangle.
Theorem~\ref{u:thm:compatible-inverse-intersection} proves this compatibility
and exhaustion.  The mixed-index route gives
\eqref{u:eq:intro-all-orders-intersection} directly.  The completed-height
identities are evaluated on the same rectangles; exhausting their heights gives
the spatial Sobolev column and then, by the equation, the same mixed
intersection.

For the velocity-time column, time Sobolev embedding applies to
\(H^m\)-valued functions: fixing \(m\) in
\eqref{u:eq:intro-all-orders-intersection} gives
\(u\in C^\infty(0,T_*;H^m)\).  Taking every finite \(m\) then gives
joint smoothness in time and space.  The pressure formula and its
differentiated recurrences determine the pressure rows from these velocity
rows.  Applying the temporal and spatial embeddings to that full
pressure-explicit jet gives joint smoothness of \((u,p)\).
Each complete representation thus has its own smoothness implication;
their compatibility identifies the derivatives used throughout the proof.

The zeroth time row of \eqref{u:eq:intro-all-orders-intersection} gives, for
each finite \(m\), \(u\in L^2(0,T_*;H^{m+2})\).  Sobolev multiplication gives
\[
 \norm{\Pp\nabla\!\cdot(u\otimes u)}_{H^m}
 \le C_m\norm u_{H^{m+2}}^2.
\]
Since \(T_*<\infty\), the nonlinear term and \(\Delta u\) belong to
\(L^1(0,T_*;H^m)\); for the linear term, explicitly,
\[
 \norm{\Delta u}_{L^1(0,T_*;H^m)}
 \le T_*^{1/2}\norm u_{L^2(0,T_*;H^{m+2})}.
\]
The equation
\[
 u_t=\nu\Delta u-\Pp\nabla\!\cdot(u\otimes u)
\]
now gives
\[
 u\in W^{1,1}(0,T_*;H^m)
 \hookrightarrow C([0,T_*];H^m)
 \qquad(m<\infty).
\]
The endpoint values agree under the Sobolev inclusions because they are limits
of one function \(u(t)\).  They define one terminal state
\[
 u_*:=u(T_*)\in\bigcap_{m<\infty}H^m.
\]
Since \(u_*\in H^\infty_\sigma\), local well-posedness restarts the equation
at \(T_*\).  Equations \eqref{u:eq:intro-pressure-recurrence}--%
\eqref{u:eq:intro-velocity-recurrence} recover the endpoint pressure and time
derivatives.  Uniqueness joins the restarted solution to \(u\), producing an
extension past \(T_*\).  This contradicts maximality and proves global
smoothness.

\clearpage
